% Fontes deste capítulo: estudo-caso/registro_conversao.md,
% estudo-caso/resultados_esperados.csv e o resultado do pytest (42 passed).

Este capítulo mostra o que aconteceu no estudo de caso. Ele descreve o que foi convertido, o resultado dos testes, os desafios observados, o que o desenvolvedor pode fazer e os limites do estudo.

\section{O que foi convertido}
\label{sec:convertido}

Neste estudo, quatro regras do SIS Eventos foram convertidas do Power Fx, a linguagem de fórmulas do Power Apps, para Python. As três primeiras foram planejadas desde o início, e a quarta foi feita a mais:

\begin{itemize}
    \item regra 01: contagem dos eventos por ano e por status, sem os registros excluídos;
    \item regra 02: valor de cada linha de custo;
    \item regra 03: total previsto de custos dos eventos aprovados de um ano;
    \item regra 04: total realizado dos eventos aprovados de um ano.
\end{itemize}

As fórmulas foram copiadas do editor do Power Apps, uma por arquivo. A conversão foi feita pelo agente Claude Code, com o modelo Claude Opus 5.5, em duas rodadas (Seção~\ref{sec:uso-agente}). O código ficou em um arquivo Python, e os testes foram escritos com a ferramenta pytest.

Os dados usados são inventados. São três tabelas, com os mesmos nomes de campos das fórmulas: eventos, linhas de custo e diárias nacionais. Os casos cobrem os anos de 2025 e 2026 e 23 linhas de custo.

Para saber o resultado certo de cada caso, as fórmulas originais foram rodadas em um aplicativo de teste do Power Apps, criado só para esse fim, com esses mesmos dados inventados. O aplicativo real não foi usado. A autora anotou o que o aplicativo mostrou em uma tabela de resultados esperados. Por exemplo, em 2026 o total previsto foi R\$~19.450,00 e o total realizado foi R\$~20.350,00. O código em Python foi comparado com essa tabela por um script (Seção~\ref{sec:testes}).

\section{Resultado da comparação e dos testes}
\label{sec:testes}

A comparação entre o código em Python e os valores esperados foi realizada por um script, executado fora da pasta do agente (Apêndice~\ref{ap:comparacao}). Dos 41 valores esperados, os 41 foram iguais aos calculados pelo código. O Quadro~\ref{quad:comparacao} apresenta o resultado por regra.

\begin{table}[htb]
\centering
\caption{Resultado da comparação com os valores esperados, por regra}
\label{quad:comparacao}
\small
\begin{tabular}{p{5.6cm}p{2.2cm}p{1.6cm}p{1.8cm}}
\toprule
\textbf{Regra} & \textbf{Valores esperados} & \textbf{Iguais} & \textbf{Diferentes} \\
\midrule
01 -- eventos por ano & 14 & 14 & 0 \\
02 -- valor da linha de custo & 23 & 23 & 0 \\
03 -- total previsto & 2 & 2 & 0 \\
04 -- total realizado & 2 & 2 & 0 \\
\midrule
Total & 41 & 41 & 0 \\
\bottomrule
\end{tabular}
\\[2pt]
\footnotesize Fonte: elaborado pela autora, com base na saída do script de comparação.
\end{table}

Nas regras 02, 03 e 04, os textos em reais gerados pelo código também foram idênticos aos da tabela de resultados esperados. Na regra 01, a lista de anos apareceu na mesma ordem da fórmula, do maior para o menor.

Além disso, os 45 testes escritos pelo agente foram aprovados (Apêndice~\ref{ap:pytest}). Eles se dividem em: 3 de leitura dos dados e formatação, 3 da regra 01, 28 da regra 02 (um por linha de custo e cinco casos especiais), 7 da regra 03 e 4 da regra 04. Nesses testes, os valores esperados estão escritos no próprio arquivo de testes, e não são lidos da tabela de resultados esperados do estudo. Por isso, eles não substituem a comparação apresentada acima.

\section{Comparação entre as duas conversões}
\label{sec:comparacao-conversoes}

As duas conversões registraram dúvidas sobre as fórmulas. A primeira as anotou no registro da conversão, e a segunda, no arquivo de dúvidas. O Quadro~\ref{quad:duvidas} apresenta os pontos que apareceram nas duas.

\begin{table}[htb]
\centering
\caption{Pontos registrados nas duas conversões}
\label{quad:duvidas}
\small
\begin{tabular}{>{\raggedright\arraybackslash}p{6.6cm}>{\raggedright\arraybackslash}p{1.4cm}>{\raggedright\arraybackslash}p{3.4cm}}
\toprule
\textbf{Ponto} & \textbf{Regras} & \textbf{Como foi registrado} \\
\midrule
A diária nacional é multiplicada pela quantidade de pessoas na regra 3, mas não na regra 2 & 2 e 3 & Nas duas conversões \\
A diária excluída é usada pela regra 2 e ignorada pela regra 3 & 2 e 3 & Nas duas conversões \\
A regra 2 não filtra as linhas excluídas & 2 & Nas duas conversões \\
A regra 3 soma todos os campos da linha, sem olhar o tipo de custo & 3 & Nas duas conversões \\
O ``-'' da função \texttt{Coalesce} nunca aparece & 4 & Nas duas conversões \\
A comparação com ``US\$'' é exata & 3 & Nas duas conversões \\
O que fazer quando o valor convertido está vazio & 2 & Nas duas conversões, com decisões diferentes \\
\bottomrule
\end{tabular}
\\[2pt]
\footnotesize Fonte: elaborado pela autora, com base no registro da primeira conversão e no arquivo de dúvidas da segunda.
\end{table}

Nos casos testados, esses pontos apareceram tanto na conversão em que o agente tinha acesso aos resultados esperados quanto na que não tinha. Isso sugere que eles vêm das fórmulas, e não da forma como o agente foi conduzido em uma das rodadas.

Houve também pontos registrados em apenas uma das conversões. Só a primeira apontou que a função \texttt{Year()} pode depender do fuso horário, que a coluna do ano se chama \texttt{Value} por causa da função \texttt{Distinct} e que não ficou claro a qual lista pertence o campo \texttt{ID} nos filtros da regra 3. Só a segunda apontou que a comparação de texto dos status é exata, que um evento sem data de início causa erro e que a ordem dos casos na descrição da regra 2 é diferente da ordem na fórmula.

O caso do valor convertido vazio mostra um limite. Quando a fórmula não diz o que fazer, cada conversão tomou uma decisão: a primeira devolveu o texto ``R\$~'', e a segunda gera um erro. Nenhuma das duas decisões pôde ser confirmada com os dados do estudo, que não têm esse caso.

\section{Desafios observados}
\label{sec:desafios-observados}

Mesmo com todos os testes passando, a leitura das fórmulas mostrou comportamentos que não estavam claros na descrição do sistema. Nos casos testados, o código em Python repete esses comportamentos, porque a ideia era ser fiel ao sistema antigo, e não corrigir as regras.

\subsection*{A regra 02 não olha o campo Excluido}

A regra 02 calcula o valor da linha mesmo quando a linha está marcada como excluída. Nos casos testados, o custo 9 está excluído e mesmo assim deu R\$~800,00. Já as regras 01, 03 e 04 deixam os registros excluídos de fora. Ou seja, a regra dos registros excluídos não está em um lugar só, e cada fórmula trata isso de um jeito. Isso se liga ao desafio das regras escondidas do Quadro~\ref{tab:desafios}, que inclui regras espalhadas em vários lugares.

\subsection*{A diária nacional não é multiplicada pelas pessoas}

Na regra 02, a diária nacional da lista é calculada como quantidade de diárias $\times$ valor da diária, sem multiplicar pela quantidade de pessoas. Nos casos testados, o custo 4 tem 2 pessoas e deu R\$~800,00, que é 2 diárias $\times$ 400. Esse comportamento só aparece quando se lê a fórmula. Por isso, esse ponto foi associado ao desafio das regras escondidas do Quadro~\ref{tab:desafios}.

\subsection*{A diária excluída ainda é usada}

A regra 02 busca o valor da diária na lista de diárias sem tirar as que estão excluídas. Nos casos testados, o custo 5 usa uma diária excluída e deu R\$~1.000,00. A regra 03, por outro lado, ignora as diárias excluídas. É o mesmo caso da regra dos excluídos espalhada em várias fórmulas, e por isso esse ponto também foi associado às regras escondidas do Quadro~\ref{tab:desafios}.

\subsection*{O ``-'' da regra 04 nunca aparece}

A regra 04 usa a função \texttt{Coalesce} para mostrar ``-'' quando não há valor. Mas o primeiro argumento sempre começa com ``R\$'', então nunca fica vazio, e o ``-'' não aparece. O código em Python também não devolve ``-''. Não foi identificada ligação direta desse ponto com os desafios do Quadro~\ref{tab:desafios}.

\subsection*{De onde vêm os valores convertidos}

Na regra 02, a diária internacional e a inscrição não fazem conta. Elas só mostram dois campos já guardados na linha: \texttt{ValorConvertido} e \texttt{InscricaoConvertida}. Neste estudo, a origem desses valores não foi identificada. Eles podem ser calculados em outra parte do sistema, que não foi vista. Isso parece com o desafio das regras escondidas do Quadro~\ref{tab:desafios}, no qual o agente não percebeu partes vizinhas do sistema. Como a origem dos campos não é conhecida, essa ligação não pôde ser confirmada.

\subsection*{A mesma conta feita de dois jeitos}

O valor das linhas de custo aparece em duas fórmulas: na regra 02, linha por linha, e na regra 03, dentro do total. As duas fazem a conta de jeitos diferentes. A regra 03 multiplica a diária nacional pelas pessoas e ignora a diária excluída, e a regra 02 não faz nenhuma das duas coisas. Por isso, em 2026, a parte do evento 1 no total da regra 03 é 6.350, enquanto a soma das linhas desse evento na regra 02, sem o custo 9, é 6.550. Isso se liga ao desafio das regras escondidas do Quadro~\ref{tab:desafios}, porque a mesma regra está em mais de um lugar.

\subsection*{O texto em reais}

O app mostra os valores como texto, no formato ``R\$~1.234,56'', que vem da função \texttt{Text()} do Power Fx. O código em Python precisou montar esse mesmo texto. Também ficou uma dúvida: quando a regra 02 não tem resultado, o app parece mostrar só ``R\$'', porque o \texttt{Text()} de um valor vazio dá texto vazio. Esse caso não foi testado. Não foi identificada ligação direta desse ponto com os desafios do Quadro~\ref{tab:desafios}.

\section{O que o desenvolvedor pode fazer}
\label{sec:desenvolvedor}

A partir do que foi observado neste estudo, sugerem-se alguns cuidados. Eles vêm de um caso só e de poucas regras, então são sugestões, e não regras para qualquer sistema.

\begin{itemize}
    \item \textbf{Listar onde cada regra aparece.} Neste estudo, a regra dos registros excluídos estava em várias fórmulas, e cada uma tratava isso de um jeito. Antes de converter, vale procurar todas as fórmulas que usam o campo \texttt{Excluido} e anotar quais olham esse campo e quais não olham.
    \item \textbf{Pedir ao agente que aponte regras duplicadas.} O valor das linhas de custo aparecia nas regras 02 e 03, com contas diferentes. Pedir ao agente que compare as fórmulas e mostre onde a mesma conta é feita de mais de um jeito pode ajudar a achar essas diferenças antes dos testes.
    \item \textbf{Perguntar à equipe a origem dos campos guardados.} Neste estudo, a origem de \texttt{Valor\allowbreak Convertido} e \texttt{Inscricao\allowbreak Convertida} não foi identificada. Perguntar ao responsável pelo sistema antes de converter evita copiar um campo sem saber como ele é calculado.
    \item \textbf{Escrever testes com centavos e casos de borda.} Nos casos testados, todos os valores eram inteiros, e vários casos ficaram sem teste. Incluir valores com centavos, anos sem eventos aprovados e linhas sem nenhum caso que se aplique ajudaria a testar o arredondamento e os comportamentos que não apareceram.
    \item \textbf{Conferir o texto na tela.} O texto em reais foi montado a partir do formato das fórmulas, e não copiado da tela. Copiar alguns textos direto do app ajudaria a confirmar o formato, o caso do ``R\$'' sozinho e o ``-'' da regra 04.
\end{itemize}

\section{Limitações do estudo}
\label{sec:limitacoes}

Este estudo tem limites que precisam ser levados em conta.

Primeiro, os dados inventados, o roteiro do aplicativo de teste e as contas à mão dos totais foram feitos por um agente de IA, o mesmo que depois realizou a primeira conversão. Na primeira conversão, esse agente teve acesso aos resultados esperados, pois o arquivo constava entre os materiais e os testes liam seus valores. Na segunda conversão, o agente recebeu apenas as fórmulas e os dados e não teve acesso a esses resultados. Mesmo assim, os dados são os mesmos, de modo que os casos de teste não são independentes do agente da primeira rodada.

Segundo, as duas conversões usaram o mesmo agente e o mesmo modelo. O agente da segunda foi instruído a trabalhar só na pasta isolada, e o histórico da sessão mostra que ele o fez, mas o isolamento não foi técnico.

Terceiro, a coluna com o texto mostrado na tela, como ``R\$~19.450,00'', não foi copiada da tela. Ela foi montada a partir do formato das fórmulas, com os valores anotados pela autora.

Quarto, todos os valores dos dados são inteiros. Por isso, nenhum caso testa centavos nem o arredondamento. Também ficaram sem teste:

\begin{itemize}
    \item um ano que só tem eventos excluídos;
    \item um ano sem eventos aprovados;
    \item uma linha de custo com campos de mais de um tipo preenchidos;
    \item uma linha da regra 02 sem nenhum caso que se aplique, ou com o valor convertido vazio.
\end{itemize}

Quinto, foram convertidas só quatro regras de um sistema maior.

Sexto, a segunda conversão levou cerca de três minutos, segundo o histórico do agente. O tempo gasto nas demais etapas, como preparar os dados e o aplicativo de teste, não foi registrado. Por isso, o custo em tempo não foi avaliado neste estudo.

Por tudo isso, o fato de o código ter coincidido com os valores esperados na primeira execução dos testes não permite concluir que a conversão é simples. Nos casos testados, o código em Python deu os mesmos resultados do aplicativo de teste. Mas os comportamentos descritos na Seção~\ref{sec:desafios-observados} só foram vistos porque as fórmulas foram lidas com atenção, e não porque algum teste falhou.

\subsection{Pontos não verificados}
\label{sec:nao-verificados}

Alguns pontos das regras ficaram sem confirmação neste estudo:

\begin{itemize}\sloppy
    \item não foi verificado com o responsável pelo sistema se a diária nacional deve ser calculada por pessoa. Na regra 02, a conta da linha não multiplica pela quantidade de pessoas, mas na regra 03 o total multiplica;
    \item não foi verificada a origem dos campos de valor convertido, \texttt{Valor\allowbreak Convertido} e \texttt{Inscricao\allowbreak Convertida};
    \item não foi verificada a origem da variável do ano (\texttt{VarAno\allowbreak Exercicio\allowbreak Selecionado}), usada nas regras 03 e 04. No aplicativo de teste, ela foi definida à mão.
\end{itemize}
